\begin{titlepage}
\centering
{\Large Notas didácticas de un especial en video\par}
\vspace{1.0cm}
{\huge\bfseries Chen Mian (陈冕), fundador de Lovart: Agent vertical, flujos de trabajo de diseño y la curva emocional del emprendimiento de aplicaciones}\par
\vspace{0.5cm}
{\Large Zhang Xiaojun Podcast\par}
\vspace{0.35cm}
{\large Elaborado a partir del video público, la transcripción hecha con GPU, la descripción del video y diagramas conceptuales propios\par}
\vspace{0.3cm}
{\large 2025-07-19\par}
\vspace{0.85cm}
\includegraphics[width=0.88\textwidth,height=0.42\textheight,keepaspectratio]{cover.jpg}\par
\vfill
\begin{tcolorbox}[width=0.92\textwidth, colback=black!2!white, colframe=black!55, sharp corners]
\textbf{Título del video}: Chen Mian, fundador de Lovart, hace un balance de estos dos años de emprendimiento de aplicaciones (video): ¡¡Qué satisfacción en este momento!! Jajajajaja\par
\textbf{Canal del video}: Zhang Xiaojun Podcast\par
\textbf{Duración del video}: 01:44:56\par
\textbf{Enlace del video}: \href{\videourl}{\nolinkurl{\videourl}}\par
\textbf{Notas sobre el material}: YouTube no tiene subtítulos disponibles; el texto se elaboró a partir de la transcripción con faster-whisper large-v3 en CUDA, la descripción del video, la portada y diagramas didácticos propios. El video es un especial de entrevista con forma de onda de audio y subtítulos, por lo que el texto no vuelve a insertar fotogramas de las personas.\par
\textbf{Página del proyecto}: \href{\repourl}{\nolinkurl{\repourl}}\par
\end{tcolorbox}
\end{titlepage}

\tableofcontents
\newpage
